\chapter{التنفيذ والاختبارات}

\begin{chaptersummary}
عرضنا في الفصل السابق تصميم النظام. ونعرض في هذا الفصل تنفيذه: التقنيات المختارة ومبرّراتها، وقواعد البيانات المنشورة، والمكتبات المعتمدة في كل طبقة، ونماذج الذكاء الصنعي التي جُرّبت وما استقرّ عليه العمل منها، وطريقة تنفيذ البثّ المباشر والسبورة والتسجيل وقاعدة المعرفة، ومعاني الإعدادات الأساسية، وبسودوكود الخوارزميتين الجوهريتين، ثمّ خطة الاختبارات ونتائجها المقاسة.
\end{chaptersummary}

\section{التنفيذ}

\subsection{التقنيات المستخدمة}

يعرض الجدول~\ref{tab:tech} التقنيات المعتمدة ومبرّر اختيار كلٍّ منها دون سواها.

\begin{table}[htbp]
  \caption{التقنيات المستخدمة في النظام ومبرّرات اختيارها.}
  \label{tab:tech}
  \centering
  \begin{tabular}{|p{3.2cm}|p{4cm}|p{7.4cm}|}
    \hline
    \textbf{التقنية} & \textbf{الغرض} & \textbf{مبرّر الاختيار} \\ \hline
    \en{.NET 10} & الخدمات الخلفية الأساسية & نظام أنواعٍ صارم يكشف صنفًا كبيرًا من الأخطاء عند الترجمة لا عند التشغيل، مع نضج \en{Entity Framework} في إدارة المعاملات والترحيلات \\ \hline
    \en{Python / FastAPI} & خدمات الذكاء الصنعي & منظومة أدوات معالجة اللغة والمستندات متوفّرة في \en{Python} وحدها عمليًّا، ويقدّم \en{FastAPI} دعمًا أصيلًا للمعالجة غير المتزامنة التي يقوم عليها خط الأنابيب \\ \hline
    \en{React 19 + Vite} & الواجهة الأمامية & دعم ناضج لواجهات \en{WebRTC} في المتصفّح، وزمن بناءٍ قصير يسمح بدورة تطويرٍ سريعة \\ \hline
    \en{PostgreSQL + pgvector} & التخزين والبحث الدلالي & يتيح البحث الشعاعي داخل القاعدة العلائقية نفسها، فلا تُضاف قاعدة بيانات شعاعية منفصلة تحتاج إلى مزامنةٍ وتشغيلٍ وصيانة \\ \hline
    \en{LiveKit} & البثّ (\en{SFU}) & إعادة توجيه انتقائية دون إعادة ترميز، مع خدمة تسجيل (\en{Egress}) مدمجة تلتقط صفحة ويب كاملة \\ \hline
    \en{RabbitMQ + MassTransit} & أحداث المجال & فكّ الارتباط الزمني بين الخدمات، ودعمٌ جاهز لنمط صندوق الصادر وإعادة المحاولة \\ \hline
    \en{MinIO / S3} & تخزين الكائنات & واجهة \en{S3} قياسية تسمح بالانتقال إلى مزوّد سحابي دون تعديل الكود، مع روابط موقّعة مؤقّتة \\ \hline
    \en{Docker Compose} & التشغيل والنشر & إعادة إنتاج البيئة كاملةً — ستّ خدمات وخمس قواعد بيانات وناقل ومخزن وخادم وسائط — بأمرٍ واحد \\ \hline
  \end{tabular}
\end{table}

\subsection{قواعد البيانات}

عُرضت مخططات الكيانات والعلاقات لكل خدمة في الفصل السابق، فلا تُعاد هنا. وإنّما يبيّن الجدول~\ref{tab:databases} ما نُشر فعلًا: خمس قواعد بيانات مستقلّة، لكلّ خدمةٍ قاعدتها.

\begin{table}[htbp]
  \caption{قواعد البيانات المنشورة والخدمة المالكة لكلٍّ منها.}
  \label{tab:databases}
  \centering
  \begin{tabular}{|p{4.2cm}|p{5.4cm}|p{5cm}|}
    \hline
    \textbf{قاعدة البيانات} & \textbf{الخدمة المالكة} & \textbf{المحرّك} \\ \hline
    \code{users\_db} & \en{UserManagementService} & \en{PostgreSQL 17} \\ \hline
    \code{classroomdb} & \en{ClassroomService} & \en{PostgreSQL} \\ \hline
    \code{streaming\_db} & \en{StreamingService} & \en{PostgreSQL 16} \\ \hline
    \code{knowledgedb} & \en{RagService} & \en{PostgreSQL + pgvector} \\ \hline
    \code{liveassistantdb} & \en{LiveAssistantService} & \en{PostgreSQL} \\ \hline
    — & \en{EmailService} & بلا قاعدة بيانات؛ مستهلِك صرف \\ \hline
  \end{tabular}
\end{table}

وتعمل كلّ قاعدةٍ في حاويةٍ مستقلّة باعتمادها الخاصّ، ولا تملك خدمةٌ سلسلة اتّصالٍ بقاعدة خدمةٍ أخرى. وامتداد \en{pgvector} مثبَّت في قاعدة خدمة المعرفة وحدها، لأنّها الوحيدة التي تخزّن أشعّة تضمين.

ويُضاف إلى القواعد العلائقية مخزنان غير علائقيين: \en{MinIO} المتوافق مع \en{S3} لتخزين الكائنات — الملفّات المرفوعة وتسجيلات الجلسات وملفّات الملخّصات — و\en{RabbitMQ} لنقل أحداث المجال بين الخدمات.

\subsection{المكتبات والأطر}

تعرض الجداول~\ref{tab:libs-dotnet} و\ref{tab:libs-python} و\ref{tab:libs-web} المكتبات المعتمدة في كل طبقة ووظيفة كلٍّ منها.

\begin{table}[htbp]
  \caption{المكتبات المعتمدة في خدمات \en{.NET}.}
  \label{tab:libs-dotnet}
  \centering
  \begin{tabular}{|p{6.2cm}|p{8.4cm}|}
    \hline
    \textbf{المكتبة} & \textbf{وظيفتها} \\ \hline
    \en{Entity Framework Core + Npgsql} & ربط الكائنات بالعلاقات، والترحيلات، وتحقيق وحدة العمل \\ \hline
    \en{MassTransit} (\en{RabbitMQ}، \en{EntityFrameworkCore}) & ناقل الأحداث وصندوق الصادر المعامل \\ \hline
    \en{SignalR} & الدفع الفوري إلى المتصفّح داخل الجلسة الحيّة \\ \hline
    \en{Livekit.Server.Sdk.Dotnet} & إنشاء الغرف، وإصدار رموز الوصول، والتحكّم بالتسجيل \\ \hline
    \en{AWSSDK.S3} & التعامل مع مخزن الكائنات \\ \hline
    \en{JwtBearer}\newline و\en{System.IdentityModel.Tokens.Jwt} & المصادقة والتحقّق من الرموز \\ \hline
    \en{Serilog} & التسجيل البنيوي للأحداث \\ \hline
    \en{AutoMapper} & التحويل بين الكيانات وبنى نقل البيانات \\ \hline
    \en{MailKit} & إرسال البريد في خدمة البريد \\ \hline
    \en{xUnit} & الاختبارات \\ \hline
  \end{tabular}
\end{table}

\begin{table}[htbp]
  \caption{المكتبات المعتمدة في خدمات \en{Python}.}
  \label{tab:libs-python}
  \centering
  \begin{tabular}{|p{5.4cm}|p{9.2cm}|}
    \hline
    \textbf{المكتبة} & \textbf{وظيفتها} \\ \hline
    \en{FastAPI + Uvicorn} & طبقة الواجهات والتشغيل غير المتزامن \\ \hline
    \en{SQLAlchemy} (غير المتزامنة) و\en{asyncpg} و\en{Alembic} & الوصول إلى قاعدة البيانات وترحيل مخططها \\ \hline
    \en{Pydantic} و\en{pydantic-settings} & التحقّق من البيانات والإعدادات المكتوبة الأنواع \\ \hline
    \en{pgvector} & عمود الأشعّة وفهرس \en{HNSW} \\ \hline
    \en{PyMuPDF} و\en{python-docx} و\en{python-pptx} و\en{Pillow} & استخراج النصّ والصور من صيغ المستندات \\ \hline
    \en{pytesseract} & التعرّف الضوئي على الحروف \\ \hline
    \en{boto3} & قراءة الملفّات المرفوعة من مخزن الكائنات \\ \hline
    \en{markdown} و\en{WeasyPrint} & تحويل الملخّص إلى ملفّ \en{PDF} \\ \hline
    \en{aio-pika} & نشر الأحداث على الناقل \\ \hline
    \en{prometheus-client} & المقاييس التشغيلية \\ \hline
    \en{pytest} و\en{pytest-asyncio} & الاختبارات \\ \hline
  \end{tabular}
\end{table}

ويجدر التنويه بأنّ حاويتَي خدمتَي \en{Python} لا تحملان أوزان نماذج ولا مكتبات استدلال — لا \en{torch} ولا \en{transformers} — فكلّ استدلالٍ نداءُ \en{HTTP} يقع خلف منفذ.

\begin{table}[htbp]
  \caption{المكتبات المعتمدة في الواجهة الأمامية.}
  \label{tab:libs-web}
  \centering
  \begin{tabular}{|p{5.4cm}|p{9.2cm}|}
    \hline
    \textbf{المكتبة} & \textbf{وظيفتها} \\ \hline
    \en{livekit-client} و\en{@livekit/components-react} & جلسة \en{WebRTC} وقناة البيانات \\ \hline
    \en{@microsoft/signalr} & استقبال الدفع الفوري من الخادم \\ \hline
    \en{TanStack Query} & إدارة حالة الخادم وذاكرتها المؤقّتة \\ \hline
    \en{Zustand} & إدارة حالة العميل \\ \hline
    \en{react-hook-form + Zod} & النماذج والتحقّق من مدخلاتها \\ \hline
    \en{React Router} & التوجيه بين الصفحات \\ \hline
    \en{Tailwind CSS 4} & التنسيق \\ \hline
    \en{i18next} و\en{react-i18next} & التوطين وتعدّد اللغات \\ \hline
    \en{react-markdown + rehype-sanitize} & عرض ناتج النموذج بعد تعقيمه \\ \hline
    \en{Vitest} و\en{Testing Library} & الاختبارات \\ \hline
  \end{tabular}
\end{table}

\subsection{نماذج الذكاء الصنعي: المجرَّب والمعتمد}

بدأ العمل بتشغيل النماذج كلّها محلّيًّا: تحويل الكلام إلى نصّ بمحرّك \en{faster-whisper} على وحدة المعالجة المركزية، والتضمين والتوليد بالنموذجين \en{qwen3-embedding} و\en{qwen2.5:7b-instruct} داخل \en{Ollama}. وكان الدافع ظاهرًا: لا كلفة تشغيل، ولا مفاتيح، ولا اعتماد على شبكةٍ قد تنقطع.

غير أنّ هذا الاختيار اصطدم بقيدٍ حاسم، هو أنّ المراحل الثلاث تعمل \emph{في آنٍ واحد} طوال الجلسة الحيّة على مضيفٍ ذي ثماني أنوية. فكان التحويل والتوليد يتنازعان الأنوية جميعها في اللحظة نفسها، حتى إنّ الردّ القصير نفسه استغرق خمس ثوانٍ مرّةً واثنتي عشرة ثانية في التالية. وقد عولج التنازع بقسمة صريحة للأنوية — خيطان للتحويل وستّة للتوليد — فعالجت القسمةُ التذبذبَ ولم تعالج البطء الأصلي.

ويبيّن الجدول~\ref{tab:stt-engines} ما جُرّب من محرّكات التحويل، والسرعة فيه منسوبةٌ إلى الزمن الحقيقي: فالقيمة دون الواحد تعني أنّ المحرّك يستهلك الصوت أبطأ ممّا تنتجه المحاضرة، ويتخلّف تخلّفًا يتراكم ولا يعوّضه تخزينٌ مؤقّت.

\begin{table}[htbp]
  \caption{محرّكات تحويل الكلام إلى نصّ التي جُرّبت، وسرعة كلٍّ منها منسوبةً إلى الزمن الحقيقي.}
  \label{tab:stt-engines}
  \centering
  \begin{tabular}{|p{4.4cm}|p{2.2cm}|p{2.4cm}|p{5.4cm}|}
    \hline
    \textbf{المحرّك} & \textbf{موضع التشغيل} & \textbf{السرعة} & \textbf{الحصيلة} \\ \hline
    \en{faster-whisper small.en} & محلّي & \en{0.2×} & أدقّ المحلّية وأبطؤها؛ لا يلاحق الكلام الحيّ \\ \hline
    \en{faster-whisper tiny.en} & محلّي & \en{1.8×} & يلاحق الكلام، لكنّه الأسوأ دقّةً \\ \hline
    \en{faster-whisper base.en} (\en{int8}، خيطان) & محلّي & أسرع من الزمن الحقيقي & التسوية المحلّية العملية، وهو البديل حين لا شبكة \\ \hline
    \en{whisper-large-v3-turbo} عبر \en{Groq} & مستضاف & نحو \en{4.5×} & الأسرع والأدقّ معًا؛ وهو المعتمد \\ \hline
    \en{gemini-flash-latest} & مستضاف & \en{0.5–0.67×} & ملاذٌ أخير عند حجب \en{Groq}؛ يهلوس كلامًا على الصمت \\ \hline
  \end{tabular}
\end{table}

فانتقل التحويل إلى المحرّك المستضاف لأنّه كان أسرع وأدقّ في آنٍ معًا، ثمّ تبعه التضمين إلى \en{gemini-embedding-001}، ثمّ تبعهما التقييم إلى \en{gemini-flash-lite}. وأفاد التضمين المستضاف فائدةً لم تكن مقصودة: فهو متعدّد اللغات، وقد قيست عتبة الانحراف نفسها صالحةً في العربية والفرنسية كما في الإنكليزية دون إعادة معايرة، بل إنّ الجملة الواحدة بالعربية ومقابلها بالإنكليزية تباعدتا \en{0.055} فقط، فالمعلّم الذي يخلط اللغتين لا يُحدِث انقسامًا كاذبًا في الأفكار.

ويعرض الجدول~\ref{tab:models} النماذج المستعملة فعليًّا في التشغيل. وقد رُوعي في اختيارها أنّ لكلّ موضعٍ قيدًا مختلفًا: فالتضمين والتحليل يقعان على المسار الحرج للجلسة الحيّة ويحكمهما الكمون، أمّا التلخيص فيجري بعد انتهاء الجلسة ويحكمه جودة الناتج.

\begin{table}[htbp]
  \caption{نماذج الذكاء الصنعي المستعملة في التشغيل وموضع كلٍّ منها.}
  \label{tab:models}
  \centering
  \begin{tabular}{|p{3.6cm}|p{5.2cm}|p{5.8cm}|}
    \hline
    \textbf{الوظيفة} & \textbf{النموذج} & \textbf{القيد الحاكم} \\ \hline
    تحويل الكلام إلى نص & \en{whisper-large-v3-turbo} عبر \en{Groq} & الكمون: يجب أن يسبق النصُّ انتهاءَ الفكرة التالية \\ \hline
    التضمين في المساعد الحيّ & \en{gemini-embedding-001} & الكمون: يقع على المسار الحرج لكشف الحدود \\ \hline
    التحليل والتدخّل & \en{gemini-flash-lite} & الكمون مع إلزام الناتج ببنية محدّدة \\ \hline
    التضمين في خدمة المعرفة & \en{gemini-embedding-001} (3072 بعدًا) & جودة الاسترجاع؛ يجري في الخلفية فلا يقيّده الكمون \\ \hline
    الإجابة والتلخيص & \en{gemini-flash} & جودة الناتج وطول السياق المتاح \\ \hline
  \end{tabular}
\end{table}

ولم يكن انتقال التضمين حرًّا كانتقال التوليد. فأشعّة نموذجين مختلفين تقع في فضاءين مختلفين، ومزجها لا يُنتج خطأً بل نتائج واثقةً خاطئة؛ ولذلك كانت الهجرة الجزئية أسوأ من عدمها. ثمّ إنّ عرض الشعاع هو نفسه عرض العمود في قاعدة البيانات وفهرسه، فاقتضى الانتقال من 1024 بعدًا إلى 3072 ترحيلًا للمخطّط وإعادةَ تضمينٍ للمقاطع المخزَّنة كلّها، وهو ما استدعى استعمال النوع \en{halfvec} كما شُرح في القسم~\ref{subsec:rag-data-model}. أمّا انتقال التوليد فحرٌّ ويمكن الرجوع عنه بلا ترحيل، إذ النصّ نصٌّ لا يُحفَظ بصيغةٍ خاصّة بنموذج.

وقد استُبقي المسار المحلّي كاملًا لا مهجورًا: يكفي متغيّرٌ واحد لإعادة أيٍّ من المراحل الثلاث إلى تشغيلها المحلّي. غير أنّه لم يُجعَل بديلًا \emph{تلقائيًّا} عند تعذّر المستضاف، وهذا اختيارٌ مقصود: التدهور الصامت إلى محرّكٍ أبطأ عشرين ضعفًا أعسر في التشخيص من عطلٍ صريح يُعلن نفسه.

وظهر بعد الانتقال عطلٌ لم يكن في الحسبان: يُجيب مزوّد التحويل بالرمز 403 عن كلّ طلب، فتموت الجلسة عند أوّل نداءٍ للتحويل. وتبيّن بالفحص أنّ المانع مانعان متراكبان لا مانعٌ واحد: الأوّل جغرافي، إذ يقع الاتصال المنزلي في بلدٍ يرفضه المزوّد ابتداءً؛ والثاني أنّ المزوّد محجوبٌ خلف شبكة توزيعٍ ترفض عناوين مراكز البيانات، وكلّ مخرج شبكةٍ خاصّة افتراضية عنوانُ مركز بيانات. فالبلد المسموح شرطٌ لازم غير كافٍ. والفارق بين المانعين لا يظهر في رمز الحالة إذ إنّه 403 في الحالتين، بل في جسم الجواب: نصٌّ عامّ يعني حجبًا عند حافّة الشبكة، وإشارةٌ صريحة إلى المفتاح تعني أنّ المفتاح هو العلّة.

ولذلك أُضيف محرّك تحويلٍ ثانٍ مستضاف من مزوّدٍ يصل حيث لا يصل الأوّل، ولم يُجعَل هو الأصل: فقياسه \en{0.5} إلى \en{0.67} من الزمن الحقيقي، أي إنّه يتخلّف عن المحاضرة بنيويًّا، كما أنّه يعيد كلامًا مختلَقًا على نوافذ الصمت لأنّه نموذج لغةٍ لا نموذج تعرّفٍ على الكلام. والخلاصة الهندسية أنّ خدمةً مستضافة ترفض المنطقة التي يجري فيها العرض أساسٌ هشّ لعملٍ يُقيَّم بالعرض المباشر، وأنّ المسار المحلّي — على بطئه — هو الضمانة الوحيدة التي لا تسقط بقرارٍ من طرفٍ ثالث.

\subsection{تنفيذ البثّ المباشر}

خادم \en{LiveKit} مُعيد توجيهٍ انتقائي (\en{SFU}): يستقبل مسارات المشاركين ويوزّعها على الآخرين دون إعادة ترميزها. وتقتصر مهمّة خدمة البثّ على إنشاء الغرفة وإصدار رمز وصولٍ لكل مشارك ثمّ الانسحاب؛ فالوسائط تنتقل بين المتصفّح وخادم \en{LiveKit} مباشرةً عبر \en{WebRTC}، دون المرور بالبوّابة العكسية ولا بالخدمة نفسها. وهذا ما يمنع الخدمة من أن تصير عنق الزجاجة عند ازدياد عدد المشاركين.

ويُحمَل دور المشارك في \emph{بيانات التعريف} (\en{Metadata}) ضمن الرمز الموقَّع، لا في اسم العرض؛ والسبب أنّ الاسم قابل للتغيير من العميل، فبناء أيّ قرارٍ تخويليّ عليه يجعل التخويل قابلًا للانتحال. ويحدّد الرمز فوق ذلك ما يُسمح لحامله بنشره من صوتٍ وصورةٍ ومشاركةِ شاشة.

وقد اعترض التنفيذَ عطلٌ في هذا الموضع: كانت كلّ جلسةٍ تنقطع بعد نحو ثلاث عشرة ثانية، ولا ينتقل صوتٌ قطّ. وأظهر سجلّ خادم الوسائط ثلاثين زوجًا من مرشَّحات \en{ICE} دون زوجٍ ناجح واحد. والسبب أنّ تفاوض \en{ICE} يقتضي تحقّق الاتجاهين معًا، وكان أحدهما وحده يعمل: فالمتصفّح يصل إلى الخادم، بينما يُعلن المتصفّح عناوين المضيف نفسه — واجهة الشبكة المحلّية، ومعها عنوان النفق حين تعمل شبكة خاصّة افتراضية — وهي عناوين غير قابلة للتوجيه من داخل الآلة الافتراضية التي تعمل فيها الحاويات. وضاعف المشكلةَ أنّ بيئة الحاويات المستعملة تمرّر حركة \en{TCP} على الواجهات كلّها بينما تمرّر \en{UDP} على العنوان المحلّي وحده؛ وهذا التفاوت هو العطل بعينه: إشارة التفاوض تصل، والوسائط لا سبيل لها أن تصل.

فاستُبدل بمحاولة إرضاء الشرط إلغاءُ الشرط نفسه: أُدخل خادم ترحيل (\en{TURN}) يخصّص المتصفّح عليه منفذًا عبر اتصالٍ صادر وتمرّ الوسائط خلاله، فلا يُنشئ خادمُ الوسائط رزمةً واحدة باتّجاه العميل، ويسقط سؤالُ أيّ العناوين قابلٌ للتوجيه من أصله. ولا يكلّف هذا شيئًا حين يعمل المسار المباشر، إذ يبقى تفاوض \en{ICE} مرجّحًا للزوج المباشر متى صحّ، فالترحيل بديلٌ عند التعذّر لا طريقٌ التفافيّ دائم. وأُلحق بذلك توحيد الإعدادات المتفرّقة التي كانت تصف المضيف في ثلاثة مواضع، في متغيّرٍ واحد يقودها جميعًا، بعد أن كان اختلافها فيما بينها مصدرًا مستقلًّا لانقطاع الجلسات.

\subsection{تنفيذ السبورة البيضاء}

للسبورة نمطا عمل: الشرح فوق الشاشة المشارَكة (\en{annotate})، والرسم على سطحٍ فارغ بنسبة \en{16:9} (\en{board}). وهي في الحالتين طبقةُ رسمٍ (\en{Canvas}) في المتصفّح، لا تُدمَج في مسار الفيديو ولا تمرّ بالخادم.

ولذلك فإنّ ما يراه الطالب ليس فيديو معدَّلًا، بل رسمًا يُعاد إنشاؤه عنده: تُرسَل نقاط المعلّم عبر \emph{قناة البيانات} في \en{LiveKit}، ويُعيد كلُّ عميلٍ رسمها على طبقته فوق الصورة. واختيرت قناة البيانات دون مركز \en{SignalR} لأنّها مفتوحة أصلًا وموثَّقة برمز الغرفة، ولأنّه لا معنى لرحلة ذهابٍ وإيابٍ إلى الخادم في ميزةٍ بصريةٍ محضة لا تُخزَّن ولا تُدقَّق؛ ولم تستلزم الميزة من ثمّ أيّ تعديلٍ في الخلفية.

وتُنقل النقاط بإحداثيات \emph{معيارية} في المجال $[0,1]$ منسوبةً إلى مستطيل \emph{محتوى} الفيديو — أي ما يتبقّى بعد الأشرطة السوداء الناتجة عن احتواء الصورة (\en{Letterboxing}) — لا بالبكسل. ولو نُقلت بالبكسل لظهرت الدائرة التي يرسمها المعلّم حول كلمةٍ ما حول كلمةٍ أخرى في كل متصفّحٍ باختلاف قياس نافذته. ويُرسل الرسم الحرّ على هيئة فروقٍ متتابعة تُجمَّع كل 50 ميلي ثانية، بينما تُرسل الأشكال الهندسية مرّةً واحدة عند رفع المؤشّر.

ولا تحفظ قناة البيانات ما مضى، فالطالب الذي ينضمّ بعد دقيقتين يجد سبورةً فارغة. ولذلك يطلب الحالة بدل انتظارها: يرسل رسالة \code{hello} عند الاتصال، فيجيبه المعلّم بالسبورة موجَّهةً إليه وحده. والطلبُ أمتنُ من أن يترقّب المعلّم القادمين، لأنّه يغطّي كذلك الحال التي يكون فيها \emph{المعلّم} هو من أعاد الاتصال فلم يشهد وصول أحد.

أمّا تجميد الصورة فيُنفَّذ بإيقاف عنصر الفيديو نفسه لا بإرسال صورةٍ ساكنة: فالعنصر المرتبط بمجرى وسائط يحتفظ بآخر إطارٍ ما دام موقوفًا ويعود إلى الحافّة الحيّة عند تشغيله، فلا يكلّف ذلك شيئًا على الشبكة. وطبقةُ الرسم لا تحجب المؤشّر إطلاقًا؛ إنّما تحجبه لوحةُ الرسم وحدها، وفي أثناء إمساك المعلّم بأداةٍ فقط، وإلّا لابتلعت النقرات الموجَّهة إلى أزرار التحكّم بالفيديو تحتها.

\subsection{تنفيذ التسجيل والتخزين}

لا تدمج خدمة \en{Egress} مسارات الوسائط، بل تفتح متصفّح \en{Chrome} بلا واجهة رسومية وتحمّل فيه صفحةً ثمّ تلتقط ما تعرضه. ولهذا أُنشئت في الواجهة الأمامية صفحة مخصّصة هي \code{/recorder}، تنضمّ إلى الغرفة مشاركًا عاديًّا وتعرض المشهد نفسه وطبقة السبورة نفسها.

وهذا هو السبب الذي يجعل رسوم السبورة تظهر في ملفّ الفيديو أصلًا: فهي مرسومة على لوحةٍ في متصفّح، ولا يمكن لقالبٍ لا يعرف إلّا المسارات أن يراها؛ أمّا هنا فتقع الرسوم في الفيديو لأنّها على الشاشة الملتقَطة. ولا ينظر أحدٌ إلى هذه الصفحة قطّ: يحمّلها آليٌّ برمزٍ في مسار العنوان، فلا تحمل ترويسةً ولا أزرار تحكّمٍ ولا مصادقة، إذ إنّ كلّ ما تعرضه يُحرَق في التسجيل حرقًا دائمًا. وهي تنتظر إشارة \code{START\_RECORDING} قبل بدء الالتقاط، ولا تُصدرها إن نقصها العنوان أو الرمز، فيفشل تسجيلٌ سيّئ الضبط فشلًا صريحًا بدل أن يسجّل ساعةً من فراغ. ويبيّن الشكل~\ref{fig:flow-03} تدفّق التسجيل كاملًا.

ويحتاج تسليم مهامّ التسجيل من خادم الوسائط إلى عامل \en{Egress} إلى وسيطٍ بينهما هو \en{Redis}. وهو ليس جزءًا من مكوّنات النظام ولا تقرؤه خدمةٌ من خدماته ولا تكتب فيه؛ وظيفته الوحيدة توزيع المهامّ داخل \en{LiveKit} نفسه، وبدونه لا تُسلَّم طلبات التسجيل أصلًا. أمّا وجهة الرفع فليست في ضبط العامل، بل تُمرّرها خدمة البثّ مع كلّ طلب: اسم الحاوية والمفتاح وعنوان المخزن. ويصل خبر الاكتمال إشعارًا موقَّعًا (\en{Webhook}) من \en{LiveKit} يُتحقَّق منه بالمفتاح المشترك. وضُبط فوق ذلك حدٌّ أقصى قدره ثلاث ساعات للتسجيل الواحد، إذ إنّ تسجيلًا لا يصل طلبُ إيقافه — لتعطّل الخدمة أو قتل الحاوية — يظلّ يعمل حتى تخلو الغرفة.

وقد ضُبطت إعدادات التسجيل بناءً على قياسٍ لا على تقدير، إذ إنّ أسوأ أنماط الفشل في التسجيل هو ملفٌّ بحجم صفر بايت. ويبيّن الجدول~\ref{tab:egress} القياسين اللذين استُند إليهما، والمقياس الحاسم فيهما هو عدد الإطارات المُسقَطة.

\begin{table}[htbp]
  \caption{قياس أثر إعدادات التسجيل في عدد الإطارات المُسقَطة خلال 284 ثانية.}
  \label{tab:egress}
  \centering
  \begin{tabular}{|p{3.6cm}|p{2.8cm}|p{2.6cm}|p{5.4cm}|}
    \hline
    \textbf{الإعداد} & \textbf{الإطارات المُسقَطة} & \textbf{نسبة التجمّد} & \textbf{النتيجة} \\ \hline
    \en{1280×720 @ 15fps} & 3708 & 67\% & غير قابل للمشاهدة \\ \hline
    \en{960×540 @ 10fps} & 0 & 28\% & انسيابيّ \\ \hline
  \end{tabular}
\end{table}

ونسبة التجمّد في السطر الثاني ليست عيبًا بل أثرٌ متوقَّع لمحتوًى ساكن: فالشريحة التي لا تتغيّر لا تُنتج إطارات جديدة. وقد خُفِّض معدّل البتّ من 4500 إلى 1200 كيلوبت في الثانية، لأنّ القيمة الافتراضية مضبوطة لمحتوًى متحرّك بدقّة عالية وكانت تكلّف نحو غيغابايت واحد في الساعة من الشرائح الساكنة. كما أُعيد ترتيب صفحة التسجيل بجعل الكاميرا في زاويةٍ صغيرة بدل شريطٍ عرضيّ كامل، فارتفعت نسبة ما تشغله المادة العلمية من الإطار من 54\% إلى 97\%.

وفُعِّل في صفحة التسجيل \en{adaptiveStream}، وهو معطَّل افتراضيًّا. فبدونه تشترك الصفحة في أعلى طبقات البثّ المتعدّد — أي شاشةٍ مشارَكة بدقّة \en{1080p} — ثمّ تصغّرها إلى إطار التسجيل في كل لقطة، داخل المتصفّح الذي هو أصلًا عنق الزجاجة. وهو لا يرفع دقّة التسجيل، فتلك محكومة بإطار الالتقاط وحده، وإنّما يحدّد أيّ طبقةٍ تُجلَب لملئه.

\subsection{تنفيذ قاعدة المعرفة}

عند رفع المعلّم ملفًّا، تحفظه خدمة الفصول في مخزن الكائنات ثمّ تُرسل طلبًا داخليًّا إلى خدمة المعرفة لتشغيل الفهرسة. ولا تنتظر الخدمةُ الطالبة اكتمال المعالجة، بل تُسجَّل المهمّة وتُعالَج في الخلفية عبر عاملٍ محدود التزامن، فلا تتعطّل عملية الرفع بانتظار استخراجٍ وتعرّفٍ ضوئيٍّ وتضمينٍ قد تستغرق دقائق. ويبيّن الشكل~\ref{fig:arch-02} مراحل هذا الخط.

وتُحتسب بصمة لمحتوى الملفّ (\en{Content Hash}) تُستعمل لجعل المعالجة خاملة التكرار (\en{Idempotent}): فإعادة إرسال الطلب نفسه لا تُعيد الفهرسة، بينما يؤدّي تغيّر المحتوى الفعلي إلى إعادتها. أمّا التعرّف الضوئي فيُطبَّق انتقائيًّا على الصفحات الممسوحة والصور الحاملة للنصّ فقط، لأنّ تطبيقه على كل صفحة يضاعف زمن المعالجة دون أن يضيف نصًّا جديدًا في الصفحات النصّية أصلًا.

وأمّا الاسترجاع فيُقصَر دائمًا على الفصل الدراسي المطلوب، لأنّ معرّف الفصل مخزَّنٌ في كيان المقطع نفسه ويدخل في الاستعلام شرطًا لا مرشِّحًا لاحقًا. وهذا ما يمنع بنيويًّا تسرّب محتوى فصلٍ إلى نتائج فصلٍ آخر.

\subsection{معاني الإعدادات الأساسية}
\label{subsec:settings}

لا تحمل إعدادات هذا النظام قيمًا اعتباطية؛ فأكثرها استقرّ على قيمته الحالية بعد قياسٍ كشف عطبًا في القيمة السابقة. وتُجمَع إعدادات خدمتَي \en{Python} في ملفّ \code{settings.py} واحد لكلٍّ منهما، ويرافق كلَّ قيمة تعليقٌ يذكر ما الذي حدث عند غيرها، حتى لا تُعاد التسوية نفسها بعد حين. وتعرض الجداول الآتية أهمّ هذه الإعدادات ومعانيها.

\begin{table}[htbp]
  \caption{إعدادات تحويل الكلام إلى نصّ ومعنى كلٍّ منها.}
  \label{tab:cfg-stt}
  \centering
  \begin{tabular}{|p{4.2cm}|p{4.4cm}|p{6cm}|}
    \hline
    \textbf{الإعداد} & \textbf{معناه} & \textbf{قيمته ومبرّرها} \\ \hline
    \code{STT\_PROVIDER} & المحرّك الذي يحوّل الكلام إلى نصّ & \en{groq} المستضاف؛ والمحلّي \en{faster\_whisper} بديلٌ لا يُنتقل إليه تلقائيًّا \\ \hline
    \code{STT\_LANGUAGE} & لغة الكلام المتوقَّعة & فارغ، أي كشفٌ تلقائي، وهو المناسب لمعلّم يخلط العربية بالمصطلح الإنكليزي \\ \hline
    \code{STT\_PAUSE\_SECONDS} & السكوت الذي يُنهي مقطعًا نصّيًّا & \en{1.5} ثانية؛ وعند \en{0.8} كانت الجمل تُقطع في منتصفها فتضيع المفردات التقنية التي يقوم عليها التقييم \\ \hline
    \code{STT\_CPU\_THREADS} & الأنوية المتاحة للتحويل المحلّي & 2 من 8، ليبقى للتوليد ما يكفي منها \\ \hline
    \code{STT\_BEAM\_SIZE} & عرض البحث عند فكّ الترميز & 5؛ وعند 1 يفوز أوّل رمزٍ محتمل فتُكتب أسماء العلم خطأً بثقة \\ \hline
    \code{STT\_INITIAL\_PROMPT} & مفردات المجال تُلقَّن للمفكِّك & فارغ؛ فقائمة المفردات جعلت المحرّك يردّد الموجّه نفسه في 11 مقطعًا من 32 \\ \hline
    \code{STT\_VAD\_FILTER} & إسقاط المقاطع غير الكلامية قبل التحويل & مفعَّل؛ يمنع هلوسة التكرار على الصمت والضجيج \\ \hline
  \end{tabular}
\end{table}

\begin{table}[htbp]
  \caption{إعدادات كشف حدود الفكرة ومعنى كلٍّ منها.}
  \label{tab:cfg-boundary}
  \centering
  \begin{tabular}{|p{4.2cm}|p{4.4cm}|p{6cm}|}
    \hline
    \textbf{الإعداد} & \textbf{معناه} & \textbf{قيمته ومبرّرها} \\ \hline
    \code{BOUNDARY\_DRIFT\_} \code{THRESHOLD} & مسافة جيب التمام التي تُعدّ عندها الجملة موضوعًا جديدًا & \en{0.25}؛ قيست أزواج الموضوع الواحد دون \en{0.159} وأزواج الموضوع الجديد فوق \en{0.305} \\ \hline
    \code{BOUNDARY\_PAUSE\_} \code{SECONDS} & السكوت الذي يُنهي الفكرة مهما كان الموضوع & 2 ثانية \emph{فوق} ما استهلكه التحويل، أي هدأة نحو \en{3.5} ثانية: ختامُ نقطةٍ لا التقاطُ نفَس \\ \hline
    \code{BOUNDARY\_MAX\_} \code{SECONDS} و\code{\_TOKENS} & الحدّان الأقصيان لطول الفكرة & 90 ثانية و400 رمز؛ شبكتا أمان تضمنان حدًّا في المونولوج الطويل \\ \hline
    \code{BOUNDARY\_MIN\_TOKENS} & أصغر فكرة يجوز أن تُطلق حدًّا & 20 رمزًا؛ شذرةٌ من كلمتين تُدمَج فيما بعدها ولا تُعدّ فكرة \\ \hline
    \code{BOUNDARY\_EMBED\_} \code{LOOKAHEAD} & عدد عمليات التضمين المسموح بتوازيها & 4؛ والقرارات تبقى مرتّبةً على كلّ حال، فالمتداخل هو الانتظار وحده \\ \hline
    \code{GEMINI\_EMBEDDING\_} \code{TASK\_TYPE} & ضبط فضاء التضمين لغرض القياس & \en{SEMANTIC\_SIMILARITY}؛ يفصل المواضيع أفضل لكنّه يضغط المسافات، فتغييره دون إعادة معايرة العتبة يعطّل الكشف صامتًا \\ \hline
  \end{tabular}
\end{table}

\begin{table}[htbp]
  \caption{إعدادات الاسترجاع وتوقيت التدخّل ومعنى كلٍّ منها.}
  \label{tab:cfg-retrieval}
  \centering
  \begin{tabular}{|p{4.2cm}|p{4.4cm}|p{6cm}|}
    \hline
    \textbf{الإعداد} & \textbf{معناه} & \textbf{قيمته ومبرّرها} \\ \hline
    \code{EMBEDDING\_DIM} & عرض شعاع التضمين & 3072؛ وهو عرض العمود في قاعدة البيانات لا إعدادًا حرًّا، ويستلزم النوع \en{halfvec} لتجاوزه حدّ فهرس \en{HNSW} \\ \hline
    \code{RETRIEVAL\_TOP\_K} & عدد المقاطع المطلوبة لكلّ فكرة & 6 \\ \hline
    \code{RETRIEVAL\_MIN\_SCORE} & الحدّ الذي تُعدّ المادة دونه غير ذات صلة & \en{0.25}؛ ودونه لا يُستدعى النموذج أصلًا، فلا تُنتَج ملاحظةٌ غير مؤصَّلة \\ \hline
    \code{FEEDBACK\_MIN\_} \code{INTERVAL\_SEC} & أقلّ زمنٍ بين تدخّلين متتاليين & 45 ثانية؛ صحّة الملاحظة وحدها لا تكفي لإرسالها \\ \hline
    \code{FEEDBACK\_} \code{CONFIDENCE\_MIN} & أقلّ ثقةٍ تُقبل بها ملاحظة & \en{0.5} \\ \hline
    \code{FEEDBACK\_DEDUP\_} \code{WINDOW\_SEC} و\code{\_SIMILARITY} & نافذة كشف التكرار وعتبة التشابه فيها & 300 ثانية و\en{0.85}؛ تمنعان إعادة الملاحظة نفسها بصياغةٍ أخرى \\ \hline
    \code{ASSISTANT\_SMOKE\_} \code{TEST} & وضع اختبارٍ يتجاوز الاسترجاع والتقييم المؤصَّل والتوقيت & \en{false} إلزامًا في التشغيل؛ فتفعيلُه يجعل المساعد يعلّق على الشرح دون مقابلته بأيّ مادة مرجعية \\ \hline
  \end{tabular}
\end{table}

ويجمع بين هذه الإعدادات أنّ بعضها مقترنٌ ببعض اقترانًا لا يظهر في اسمه: فعتبة الانحراف مقيسةٌ لمزوّد تضمينٍ بعينه ولنمط مهمّةٍ بعينه، وعرض الشعاع هو عرض العمود في قاعدة البيانات، وعدد أنوية التحويل مقابلٌ لعدد أنوية التوليد على المضيف نفسه. وتغييرُ طرفٍ من هذه الأزواج دون طرفه الآخر لا يُنتج خطأً ظاهرًا بل سلوكًا خاطئًا صامتًا، وهو ما يفسّر إسهاب التعليقات المرافقة لها في الكود.

\subsection{بسودوكود الخوارزميتين الجوهريتين}

نعرض فيما يأتي الخوارزميتين اللتين يقوم عليهما المساعد الحيّ. وقد كُتبتا بالإنكليزية لأنّ معرّفات الكود إنكليزية أصلًا، ولأنّ خلط اتّجاهي الكتابة في كتلةٍ برمجية يشوّه ترتيب رموزها.

\subsubsection{كشف حدود الفكرة التعليمية}

قام كشف الحدود في أوّل تنفيذه على السكوت وحده، وفشل الأساس فشلًا بنيويًّا لا يُعالَج بالمعايرة: فمحرّك التحويل يُنهي المقطع عند كلّ سكتة تنفّس، فلو عُدّت كلّ سكتةٍ حدًّا لصارت كلّ جملةٍ فكرةً مستقلّة، ولضاع المعنى الذي يُفترض أن يُقيَّم. فاستُبدل بالمقياس الزمني مقياسٌ دلالي: يُقاس الانحراف الدلالي (\en{Semantic Drift}) بين شعاع الفكرة الجارية وشعاع المقطع الجديد، وتُتَّخذ السكتات الطويلة والحدود القصوى شبكات أمانٍ فحسب.

\begin{english}
\begin{lstlisting}[caption={},captionpos=b]
INPUT: stream of finalized transcript segments
STATE: buffer of segments, running mean vector, token count

for each finalized segment S (interim text is discarded):
    V <- embed(S.text)                    # prefetched, see below
    gap <- S.start_ms - previous_segment.end_ms

    # (A) a long silence BEFORE this segment closes the running idea
    if buffer is emittable and gap >= BOUNDARY_PAUSE:
        emit close(trigger = PAUSE)

    # (B) place this segment
    representative <- mean(buffer.vectors)
    if buffer is emittable and cosine_distance(representative, V) >= DRIFT:
        completed <- close(trigger = DRIFT)   # topic changed
        start_new_buffer(S, V)
        emit completed
    else:
        append(S, V) to buffer

    # (C) unconditional safety nets
    if buffer.duration >= MAX_SECONDS:  emit close(trigger = TIME_CAP)
    if buffer.tokens   >= MAX_TOKENS:   emit close(trigger = TOKEN_CAP)

at end of stream:
    if buffer is emittable: emit close(trigger = PAUSE)

# "emittable" means the buffer holds at least MIN_TOKENS: a two-word
# fragment must never trigger drift, it merges forward into the next idea.
\end{lstlisting}
\end{english}

وقد ضُبطت العتبة بقياسٍ فرّق بين أزواج الموضوع الواحد وأزواج الموضوعين، ووُضعت بينهما بهامشٍ من الطرفين مع ميلٍ متعمّد نحو الدمج، لأنّ شقّ الشرح الواحد شطرين أسوأ من حمل جملةٍ من الموضوع التالي مع سابقه.

ويُنفَّذ التضمين هنا باستباقٍ محدود (\en{Bounded Prefetch}): إذ تُطلَق عدّة عمليات تضمين متوازية، غير أنّ \emph{القرارات} تبقى مرتّبة ترتيبًا صارمًا. والسبب أنّ شعاع المقطع لا يعتمد على سابقه، أمّا كون المقطع يبدأ فكرةً جديدة فيعتمد على الحالة التي تركها ما قبله. ولولا هذا الاستباق لتناوبت المرحلتان: لا يُحوَّل كلامٌ إلى نصّ ما دام تضمينٌ جاريًا؛ وقد قيس أنّ التضمين المتتابع استغرق وحده \en{90.6} ثانية داخل محاضرةٍ طولها 143 ثانية، أي إشغالٌ يبلغ 61\% من الزمن المتاح. وحُدَّ التوازي لأنّ المضمِّن خدمةٌ مستضافة محدودة المعدّل، فالانتشار غير المحدود يحوّل دفقة كلامٍ إلى دفقة رفضٍ من الخادم.

\subsubsection{التقييم مقابل المادة المرجعية وتوقيت التدخّل}

\begin{english}
\begin{lstlisting}[caption={},captionpos=b]
on each COMPLETED IDEA:
    chunks   <- rag_service.retrieve(classroom_id, idea.text, TOP_K)
    relevant <- [c for c in chunks if c.score >= MIN_SCORE]

    if relevant is empty:
        return NO_FEEDBACK          # the model is NOT called at all
                                    # -> no ungrounded claim can be produced

    outcome <- brain.evaluate(idea, relevant)   # constrained JSON output

    if outcome has no suggestion:
        return NO_FEEDBACK

    # pacing: holding a correct remark is not enough to fire it
    if now - last_delivery < MIN_INTERVAL:
        return SUPPRESSED(reason = TOO_SOON)

    deliver_privately(teacher, outcome.suggestion)
    last_delivery <- now

# any exception in retrieval or evaluation degrades to NO_FEEDBACK:
# a failed evaluation must never break the live loop.
\end{lstlisting}
\end{english}

ويستحقّ الشرطان الآتيان الوقوف عندهما. أوّلهما أنّ النموذج \emph{لا يُستدعى إطلاقًا} حين لا يعود الاسترجاع بمادّةٍ ذات صلة؛ فهذا يقطع الطريق بنيويًّا على أيّ ملاحظةٍ غير مؤصَّلة في مادة الفصل، إذ لا يمكن للنموذج أن يدلي برأيٍ لم يُسأل عنه. وثانيهما مُوقِّت التدخّل: فالمساعد الذي يقاطع المعلّم كلّما ملك ملاحظةً صحيحة يتحوّل من عونٍ إلى مصدر تشويش، ولذلك يُفرَض حدٌّ أدنى زمنيّ بين التدخّلات مهما كانت الملاحظة وجيهة.

\section{خطة الاختبارات ونتائجها}
\label{sec:test-plan}

\subsection{منهجية الاختبار}

أعسر ما واجه التنفيذ لم يكن ما توقّف عن العمل، بل ما ظلّ يعمل وهو خاطئ: راية ضبطٍ تجاوزت الاسترجاع والتقييم المؤصَّل فجعلت المساعد يعلّق تعليقًا معقولًا دون مقابلته بأيّ مادة مرجعية، ومحرّك تحويلٍ ردّ قائمة المفردات التي لُقِّنها نصًّا في 11 مقطعًا من 32، وملخّصٌ وُلِّد ورُفِع ثمّ ضاع خبره بخطأ تسميةٍ في عنوان الناقل. ولم يُسقط أيٌّ من هذه الأعطال النظامَ ولم يُظهر رسالة خطأ، ولم يُكشَف واحدٌ منها إلّا بقياسٍ أو بقراءة سجلّ. ولذلك انصبّ التركيز في خطة الاختبار على المواضع التي يكون الخطأ فيها صامتًا.

واعتُمد في الاختبار مبدآن. الأوّل أنّ كل تبعيةٍ خارجية تقع خلف واجهة، ويقابلها في الاختبار بديلٌ مكتوب باليد (\mbox{\en{Hand-written Fake}}) لا كائنٌ مولَّد بإطار محاكاة (\en{Mocking Framework}). والثاني أنّ مجموعة الاختبارات يجب أن تعمل دون قاعدة بيانات ولا شبكة ولا نموذجٍ ولا حاوية.

وثمرة المبدأين أنّ المجموعة كاملةً — على كثرة اختباراتها — تُنجَز في أقلّ من دقيقة، فتصير تشغيلًا اعتياديًّا بعد كل تعديل لا إجراءً يُؤجَّل إلى ما قبل التسليم.

\subsection{النتائج المقاسة}

يعرض الجدول~\ref{tab:tests} نتيجة تشغيل المجموعات الستّ كاملةً بتاريخ 3 آب 2026.

\begin{table}[htbp]
  \caption{نتائج تشغيل مجموعات الاختبار الستّ كاملةً بتاريخ 3 آب 2026.}
  \label{tab:tests}
  \centering
  \begin{tabular}{|p{4.6cm}|p{1.9cm}|p{1.5cm}|p{1.5cm}|p{1.5cm}|p{1.5cm}|}
    \hline
    \textbf{المجموعة} & \textbf{الأداة} & \textbf{ناجح} & \textbf{مؤجَّل} & \textbf{فاشل} & \textbf{الزمن} \\ \hline
    \en{UserManagementService} & \en{xUnit} & 120 & 0 & 0 & \en{0.5 s} \\ \hline
    \en{ClassroomService} & \en{xUnit} & 296 & 0 & 0 & \en{4 s} \\ \hline
    \en{StreamingService} & \en{xUnit} & 104 & 0 & 0 & \en{3 s} \\ \hline
    \en{RagService} & \en{pytest} & 235 & 9 & 0 & \en{12.6 s} \\ \hline
    \en{LiveAssistantService} & \en{pytest} & 299 & 3 & 0 & \en{4.2 s} \\ \hline
    الواجهة الأمامية & \en{Vitest} & 219 & 0 & 0 & \en{32 s} \\ \hline
    \textbf{المجموع} & & \textbf{1273} & \textbf{12} & \textbf{0} & \\ \hline
  \end{tabular}
\end{table}

فمجموع ما كُتب 1285 اختبارًا، نجح منها 1273 ولم يفشل أيٌّ منها. أمّا الاثنا عشر المؤجَّلة فتتخطّى نفسها بشرطٍ صريح لا تخفق: سبعةٌ منها تحتاج إلى تنصيب محرّك التعرّف الضوئي على جهاز التشغيل، واثنان يحتاجان إلى قاعدة بيانات حيّة عبر المتغيّرين \code{TEST\_DATABASE\_URL} و\code{TRANSCRIPT\_TEST\_DB\_URL}، واثنان يحتاجان إلى ملفّ صوتي مرجعي لاختبار تحويل الكلام إلى نصّ بنموذجٍ حقيقي، وواحدٌ يختبر التعافي من فشل التعرّف الضوئي. وقد فُضِّل تخطّيها المشروط على حذفها، ليبقى تشغيلها ممكنًا في بيئةٍ مهيّأة دون تعديل الكود.

\subsection{أمثلة على ما تغطّيه الاختبارات}

اختير التركيز على المواضع التي يكون الخطأ فيها \emph{صامتًا}، أي لا يُسقط النظام بل يُنتج نتيجةً خاطئة تبدو سليمة:

\begin{enumerate}
  \item \textbf{هندسة السبورة}: تحويل الإحداثيات المعيارية إلى مستطيل المحتوى بعد الاحتواء. وخطأٌ هنا يضع الرسم في موضعٍ قريبٍ من الصحيح، فلا يُلاحَظ إلا بالمقارنة الدقيقة.
  \item \textbf{بروتوكول السبورة}: ترميز الرسائل وفكّها، ورفض غير الصالح منها، وتقسيم ما يتجاوز الحدّ الأقصى لحزمة القناة الموثوقة.
  \item \textbf{كشف حدود الأفكار}: سلوك الخوارزمية عند الانحراف والسكوت والحدود القصوى، وعند المقاطع القصيرة التي يجب ألّا تُطلق حدًّا. ويُستعمل فيها مضمِّنٌ بديل حتميّ الناتج، فلا يلزم نموذجٌ حقيقي.
  \item \textbf{عزل مفتاح الإجابة}: التحقّق من أنّ ما يُرسَل إلى الطالب لا يمكن أن يحمل حقل الإجابة الصحيحة، إذ لا تملك بنية نقل البيانات الموجَّهة إليه هذا الحقل أصلًا، ويقع التصحيح كلّه على الخادم.
  \item \textbf{اتّساق أبعاد التضمين}: رفض الإقلاع مبكرًا إذا خالف عرض الشعاع الإعداد \code{EMBEDDING\_DIM}، بدل ظهور الخطأ لاحقًا عند الإدراج.
  \item \textbf{سياسة قطع الاتصال}: قرار إعادة الوصل أو الخروج من الجلسة، وهو منطقٌ يصعب اختباره يدويًّا لأنّه يقتضي افتعال انقطاعٍ حقيقي.
\end{enumerate}

\subsection{الفجوات المعروفة وخطة معالجتها}

يقتصر المنفَّذ حاليًّا على اختبارات الوحدة والمكوّنات. ويعرض الجدول~\ref{tab:testgaps} ما لم يُنفَّذ بعدُ صراحةً، مع الخطة المقترحة لكلٍّ منه.

\begin{table}[htbp]
  \caption{الفجوات القائمة في خطة الاختبار والخطة المقترحة لمعالجتها.}
  \label{tab:testgaps}
  \centering
  \begin{tabular}{|p{3.2cm}|p{4.2cm}|p{7.2cm}|}
    \hline
    \textbf{نوع الاختبار} & \textbf{الحالة} & \textbf{الخطة} \\ \hline
    اختبارات التكامل & مشروعٌ مُهيَّأ في خدمة إدارة المستخدمين لكنّه خالٍ من أيّ اختبار & تشغيل الخدمة مقابل قاعدة بيانات حقيقية عبر \en{Testcontainers}، وتغطية مسار التسجيل والموافقة والدخول من طرفٍ إلى طرف \\ \hline
    اختبارات الأمان & غير منفَّذة & التحقّق من رفض المسارات الداخلية بلا السرّ المشترك، ومن منع الوصول العابر للأدوار، ومن انتهاء صلاحية الرموز، ومن عدم إعادة استعمال تحدّي العامل الثاني \\ \hline
    اختبارات الحمل & غير منفَّذة & قياس سعة الجلسة الحيّة بعدد مشاركين متزايد، ورصد إسقاط الإطارات في التسجيل تحت الضغط، وقياس كمون خط أنابيب المساعد الحيّ \\ \hline
  \end{tabular}
\end{table}

ويجدر التنبيه إلى أنّ اختبارات الأمان المذكورة هي اختبارات \emph{آلية}؛ أمّا الضوابط نفسها — من عزل المسارات الداخلية وحراستها بسرٍّ مشترك، وفصل بنيتي نقل البيانات للاختبارات، وتخزين بصمات الرموز الحسّاسة لا نصوصها — فمنفَّذةٌ في الكود وموصوفةٌ في هذا الفصل وفي الفصل السابق. والناقص هو ما يحرسها من الانكسار في تعديلٍ لاحق، لا الضوابط ذاتها.

\transition{وبهذا يكتمل عرض النظام تصميمًا وتنفيذًا، فننتقل إلى خاتمة التقرير لتلخيص ما أُنجز وبيان آفاق تطويره.}
